\needspace{4\baselineskip}
\section{Evaluation Setup}\label{sec:eval-env-details}

\subsection{Environment and Invocation}

\paragraph{Harness version selection.} For each model we report the earliest publicly available harness version that supported that model and could execute the benchmark end-to-end. For older models whose launch-era harness was unavailable or incompatible, we used the nearest later compatible version.

\paragraph{Environment.} The container image installs all languages required by the problem set alongside a shared tooling baseline. We derived this baseline by identifying commands whose absence caused failures across all harnesses; commands that failed on only one harness were excluded to avoid biasing the environment toward a particular agent. Installed packages, shell history, and agent session data reset between checkpoints. The benchmark problems are language-agnostic by design, but current experiments evaluate only the Python track.

\paragraph{Invocation.} Following Terminal Bench~\citep{tb2}, we install Claude Code~\citep{claude_code} and Codex~\citep{codex} directly and invoke each in headless mode. Specific versions appear in \autoref{tab:agent-versions}.

\paragraph{Reasoning effort.} For Codex we set the reasoning effort parameter to \texttt{high}. For Claude Code we configure the thinking-token budget via the environment variable following Anthropic's published mapping.

\subsection{Agent Harness Versions}\label{sec:agent-versions}

\begin{table}[h]
\centering
\caption{Harness version and reasoning effort used for each model in the main evaluation.}
\label{tab:agent-versions}
\small
\begin{tabular}{lllc}
\toprule
Model & Harness & Version & Reasoning \\
\midrule
Opus~4.5      & Claude Code & 2.0.51  & high \\
Opus~4.6      & Claude Code & 2.1.32  & high \\
Sonnet~4.6    & Claude Code & 2.1.44  & high \\
\midrule
GPT~5.2 Codex & Codex CLI & 0.93.0  & high \\
GPT~5.3 Codex & Codex CLI & 0.98.0  & high \\
GPT~5.3 Spark & Codex CLI & 0.100.0 & high \\
GPT~5.4       & Codex CLI & 0.110.0 & high \\
GPT~5.4 Mini  & Codex CLI & 0.110.0 & high \\
\midrule
Composer~2    & Cursor CLI  & 2026.04.13-a9d7fb5 & none \\
GLM~5.1       & Claude Code & 2.1.44  & high \\
Kimi~K2.5     & Claude Code & 2.1.44  & high \\
Kimi~K2.6     & Kimi CLI    & 1.37.0  & high \\
MiniMax~M2.7  & Claude Code & 2.1.44  & high \\
\bottomrule
\end{tabular}
\end{table}

\subsection{\texttt{just-solve} (Baseline)}\label{sec:prompt-intervention-raw}

The minimal baseline used for all primary evaluations. Each agent receives a Jinja template as its system prompt. The \texttt{is\_continuation} flag is false for the first checkpoint of a problem and true for all subsequent checkpoints. The checkpoint specification is injected verbatim via \texttt{spec}.

\begin{lstlisting}[style=specblock, caption={\texttt{just-solve} prompt template}, label={lst:just-solve-prompt}]
Implement a program that 100% solves the specification.
That is all you need to do.
{% if not is_continuation -%}
Use a virtual environment and ensure that a
`requirements.txt` is present with any dependencies
you need to solve the problem.
{% else -%}
Keep using the same virtual environment you started with,
update `requirements.txt` with any new dependencies
you need.
{% endif -%}

Your task is:
{{ spec.strip() }}
\end{lstlisting}

\subsection{\texttt{anti\_slop}}

Explicitly instructs the agent to avoid verbose patterns, defensive over-engineering, and unnecessary abstractions.

\begin{lstlisting}[style=specblock, caption={\texttt{anti\_slop} prompt template}, label={lst:anti-slop-prompt}]
{%- if not is_continuation -%}
You are an exceptional python software engineer and you
to need to implement a spec. Your instructions are:
- Write the python script that satisfies the spec
  completely.
- Use a virtual environment and ensure that a
  `requirements.txt` is present with any dependencies
  you need to solve the problem.
{%- else -%}
You are an exceptional python software engineer and you
to need to implement a spec. You are updating your code
to match an extension of the spec. Here are your
instructions:
- Keep using the same virtual environment you started
  with, update `requirements.txt` with any new
  dependencies you need.
- Focus only on adding in the new features/changes below.
- Make sure you test any examples provided in the task
  description
{%- endif %}
- You ONLY work in this directory.
- Follow best coding practices:
  - Group functions into files based on related
    functionality
  - Keep your code clean
- No god functions/classes.
- Make sure the code is documented appropriately so that
  it is easy to pick up.
- Minimize the following gotchas:
  - Extra defensive checks or try/catch blocks that are
    abnormal.
  - Casts to get around type checking
  - Variables that are only used a single time after
    declaration.
  - Extra comments that a human wouldn't add.
  - Trivial wrappers
  - Heavy nesting
  - If/Else ladders
  - A ton of helper methods

Your task is:
{{ spec.strip() }}
\end{lstlisting}

\subsection{\texttt{plan\_first}}

Requires the agent to plan its approach before writing code.

\begin{lstlisting}[style=specblock, caption={\texttt{plan\_first} prompt template}, label={lst:plan-first-prompt}]
You are an expert programmer and need to implement a task.
{% if not is_continuation -%}
Use a virtual environment and ensure that a
`requirements.txt` is present with any dependencies
you need to solve the problem.
{% else -%}
Keep using the same virtual environment you started with,
update `requirements.txt` with any new dependencies
you need.
{% endif -%}

Here are the steps you should always follow:
1. Before coding plan out what you need to implement.
2. Write the simple solution first.
3. Ensure it is 100% correct and you have covered all
   edge cases.
4. Refactor to ensure the code is high quality.

Here are the basic style rules you must follow:
- Make sure the code is documented appropriately so that
  it is easy to pick up.
- Minimize the following gotchas:
  - Extra defensive checks or try/catch blocks that are
    abnormal.
  - Casts to get around type checking
  - Variables that are only used a single time after
    declaration.
  - Extra comments that a human wouldn't add.
  - Trivial wrappers
  - Heavy nesting
  - If/Else ladders
  - A ton of helper methods
- Follow best coding practices:
  - Group functions into files based on related
    functionality
  - Keep your code clean

Your task is:
{{ spec.strip() }}
\end{lstlisting}
